% ============================================================================
% AVALIAÇÃO GERAL INTEGRADA — SIMULADO FINAL COMPLETO (CHF662 — DCHF/UEFS)
% ============================================================================
\clearpage
\phantomsection
\addcontentsline{toc}{chapter}{Avaliação Geral Integrada — Simulado Final Completo (CHF662)}
\chapter*{Avaliação Geral Integrada — Simulado Final Completo}
\markboth{AVALIAÇÃO GERAL INTEGRADA}{AVALIAÇÃO GERAL INTEGRADA}

\begin{tcolorbox}[
  breakable, enhanced,
  colback=bgSoft, colframe=uefsNavy,
  borderline west={4pt}{0pt}{uefsGold},
  boxrule=0.8pt, arc=3pt,
  title={\bfseries CABEÇALHO DE SIMULADO PRESENCIAL — CHF662 ANÁLISE AMBIENTAL (DCHF/UEFS)},
  coltitle=white, colbacktitle=uefsNavy,
  left=10pt, right=10pt, top=7pt, bottom=7pt
]
\small
\begin{tabularx}{\textwidth}{@{}X p{4.2cm}@{}}
  \textbf{Discente:} \hrulefill & \textbf{Matrícula:} \hrulefill \\[5pt]
  \textbf{Curso:} Geografia (DCHF/UEFS) \quad\textbf{Semestre:} 2026/2 & \textbf{Data:} \underline{\hspace{0.9cm}}/\underline{\hspace{0.9cm}}/2026 \\[5pt]
  \textbf{Tempo sugerido:} 3h00 (180 min) & \textbf{Pontuação:} 10,0 pts (Obj.: 6,0 $|$ Diss.: 4,0) \\
\end{tabularx}
\vspace{4pt}

{\footnotesize\color{uefsMuted}\textbf{Instruções gerais:} Este simulado abrange os \textbf{5 Módulos} e os \textbf{15 Capítulos} da disciplina CHF662. A \textbf{Parte I} contém 15 questões objetivas (0,4 ponto cada; apenas uma alternativa correta por questão), no padrão do ENADE e de concursos públicos para Geógrafo e Analista Ambiental (IBAMA, INEMA e órgãos estaduais). A \textbf{Parte II} contém 5 questões dissertativas integradoras (0,8 ponto cada), no padrão da prova presencial da UEFS. Realize o exame sem consulta ao gabarito, marque suas respostas na Folha de Respostas imprimível e, somente ao final, utilize o Gabarito Comentado, o Espelho de Correção e o Painel de Diagnóstico Final.}
\end{tcolorbox}

% ============================================================================
\section*{Parte I — Questões Objetivas Integradoras (Módulos 1 a 5)}
\addcontentsline{toc}{section}{Parte I — Questões Objetivas Integradoras (Q01–Q15)}
% ============================================================================

\subsection*{Módulo 1 — Fundamentos Epistemológicos e Conceituais (Capítulos 1 a 4)}

\noindent\textbf{Questão 01} {\footnotesize\color{uefsMuted}[Inédita autoral — estilo ENADE / Geografia $\bullet$ Caps. 1 e 2]}\par
A crise ambiental contemporânea levou autores como Enrique Leff (2001) e Carlos Walter Porto-Gonçalves (2006) a questionarem o modelo clássico de ciência moderna herdado da Revolução Industrial. Em uma audiência pública sobre a implantação de um polo minerário no Semiárido baiano, um parecer técnico afirmou que ``os conflitos pela água decorrem unicamente da escassez climática natural, cabendo apenas à engenharia hidráulica resolver o problema de forma neutra''. À luz da epistemologia ambiental e da análise crítica da relação sociedade--natureza, essa afirmação é metodologicamente inadequada porque:
\begin{alternativas}
  \item desconsidera que a natureza possui leis físicas imutáveis, razão pela qual o planejamento ambiental deve restringir-se exclusivamente às variáveis climáticas e hidrológicas.
  \item reduz a questão ambiental a um conflito de interesses puramente jurídicos, omitindo a obrigatoriedade de obras estruturais de transposição de bacias hidrográficas.
  \item naturaliza problemas socioambientais historicamente produzidos e ignora que a apropriação desigual dos recursos hídricos envolve racionalidade econômica, poder territorial e saberes plurais.
  \item contraria a separação necessária entre Geografia Física e Geografia Humana, que impede o geógrafo de opinar sobre obras hidráulicas e projetos econômicos.
  \item superestima a influência das comunidades locais sobre o regime de chuvas do Semiárido, atribuindo à ação humana variações que dependem apenas da circulação atmosférica.
\end{alternativas}

\vspace{6pt}
\noindent\textbf{Questão 02} {\footnotesize\color{uefsMuted}[Inédita autoral — estilo Concurso Analista Ambiental / INEMA $\bullet$ Cap. 3]}\par
Na análise integrada da paisagem proposta por Georges Bertrand (2004), o sistema tripolar \textbf{GTP} articula três entradas conceituais complementares para superar a fragmentação entre o meio físico-biótico e a dinâmica social, dialogando diretamente com a Geografia Socioambiental discutida por Francisco Mendonça (2001). Assinale a alternativa que associa \textbf{corretamente} cada polo do sistema GTP à sua dimensão analítica predominante.
\begin{alternativas}
  \item \textbf{Geossistema}: dinâmica natural e funcionamento físico-ecológico; \textbf{Território}: apropriação, gestão, recurso econômico e relações de poder; \textbf{Paisagem}: dimensão cultural, perceptiva e identitária do espaço.
  \item \textbf{Geossistema}: apenas o relevo e a geologia estrutural; \textbf{Território}: cobertura vegetal nativa sem uso humano; \textbf{Paisagem}: exclusivamente a representação cartográfica em escala de detalhe.
  \item \textbf{Geossistema}: zoneamento jurídico instituído por lei; \textbf{Território}: bacia hidrográfica fechada sem interferência antrópica; \textbf{Paisagem}: soma aleatória de elementos abióticos independentes.
  \item \textbf{Geossistema}: percepção estética da comunidade local; \textbf{Território}: equilíbrio químico dos horizontes pedológicos; \textbf{Paisagem}: infraestrutura rodoviária e malha urbana.
  \item \textbf{Geossistema}: substituto definitivo dos conceitos de território e região; \textbf{Território}: unidade estritamente biológica; \textbf{Paisagem}: sinônimo técnico de bioma na escala continental.
\end{alternativas}

\vspace{6pt}
\noindent\textbf{Questão 03} {\footnotesize\color{uefsMuted}[Inédita autoral — estilo Concurso IBAMA / Analista Ambiental $\bullet$ Cap. 4]}\par
Em um diagnóstico preliminar de riscos geomorfológicos em encostas periurbanas, uma equipe aplicou a distinção conceitual de Sánchez (2013) entre \textit{atividade}, \textit{aspecto ambiental} e \textit{impacto ambiental}, bem como a matriz conceitual de risco $R = P \times V$, na qual o Perigo ($P$) e a Vulnerabilidade ($V$) variam de $1$ (baixo) a $4$ (muito alto/crítico), classificando o risco como Baixo ($1$ a $3$), Médio ($4$ a $6$), Alto ($8$ a $9$) e Crítico ($12$ a $16$). Foram avaliados três setores ocupados:
\begin{itemize}[leftmargin=1.8em, itemsep=1.5pt, topsep=3pt]
  \item \textbf{Setor A:} corte irregular de talude com $P = 4$ e $V = 3$;
  \item \textbf{Setor B:} encosta de média declividade com drenagem parcial, $P = 3$ e $V = 2$;
  \item \textbf{Setor C:} planície marginal sujeita a alagamento sazonal, $P = 2$ e $V = 4$.
\end{itemize}
Suponha que um programa de drenagem, estabilização e reassentamento preventivo reduza a vulnerabilidade do \textbf{Setor A} de $V = 3$ para $V = 1$, mantendo-se a ameaça física natural ($P = 4$). Considerando esses dados e a cadeia causal da análise de impactos, é \textbf{correto} afirmar que:
\begin{alternativas}
  \item antes da intervenção, o Setor C apresentava risco superior ao do Setor A, pois a vulnerabilidade tem peso quadrático na equação de risco.
  \item o corte irregular do talude é o impacto ambiental, enquanto o deslizamento sobre moradias corresponde ao aspecto ambiental da ocupação.
  \item mesmo após a redução da vulnerabilidade do Setor A para $V = 1$, o risco permanecerá na classe Alto ($R = 8$), exigindo a interdição total da bacia.
  \item inicialmente o Setor A possui risco Crítico ($R = 12$), seguido pelo Setor C ($R = 8$, Alto) e pelo Setor B ($R = 6$, Médio); após a intervenção no Setor A, seu risco cai para $R = 4$ (Médio), redução relativa de $66,7\%$.
  \item a exposição do solo às chuvas intensas e o assoreamento do córrego a jusante são ambos classificados como atividades humanas primárias na cadeia causal.
\end{alternativas}

% ----------------------------------------------------------------------------
\subsection*{Módulo 2 — AIA e Análise Ambiental Integrada (Capítulos 5 e 6)}

\noindent\textbf{Questão 04} {\footnotesize\color{uefsMuted}[Inédita autoral — estilo Concurso IBAMA / Analista Ambiental $\bullet$ Cap. 5]}\par
Na Avaliação de Impacto Ambiental (AIA), conforme Sánchez (2013) e as diretrizes da Resolução CONAMA nº 001/1986, a definição das medidas de controle ambiental deve obedecer rigorosamente à \textbf{hierarquia de mitigação}. Durante o licenciamento de uma linha de transmissão que inicialmente cruzaria um remanescente contínuo de mata ciliar e uma vereda, a equipe técnica apresentou quatro propostas distintas para análise do órgão ambiental:
\begin{enumerate}[label=\Roman*., leftmargin=2.0em, itemsep=1.5pt, topsep=3pt]
  \item alterar o traçado da linha para desviar completamente da vereda e da mata ciliar;
  \item executar o apoio à criação de uma Unidade de Conservação equivalente pelos impactos residuais inevitáveis;
  \item implantar técnicas de lançamento de cabos por drones e controle rigoroso de drenagem nas praças de montagem para reduzir a área suprimida;
  \item recuperar a cobertura vegetal nas áreas temporariamente afetadas pela instalação das torres por meio de PRAD e bioengenharia.
\end{enumerate}
A sequência que ordena essas medidas da \textbf{maior para a menor prioridade} na hierarquia de mitigação (Evitar $\to$ Minimizar $\to$ Recuperar/Reabilitar $\to$ Compensar) é:
\begin{alternativas}
  \item III -- I -- IV -- II.
  \item I -- III -- IV -- II.
  \item I -- IV -- III -- II.
  \item II -- IV -- III -- I.
  \item IV -- I -- III -- II.
\end{alternativas}

\vspace{6pt}
\noindent\textbf{Questão 05} {\footnotesize\color{uefsMuted}[Inédita autoral — estilo Concurso Órgão Ambiental Estadual $\bullet$ Cap. 5]}\par
Para reduzir a subjetividade na etapa de valoração de impactos de um Estudo de Impacto Ambiental (EIA), adotou-se uma matriz ponderada em que o Índice de Significância ($I_S$) de cada impacto negativo é calculado por $I_S = \sum_{j=1}^{4} w_j \cdot n_j$, sendo as notas $n_j$ atribuídas de $1$ (menor severidade) a $4$ (maior severidade) nos critérios: \textbf{Magnitude} ($w_1 = 0,35$), \textbf{Abrangência espacial} ($w_2 = 0,25$), \textbf{Duração} ($w_3 = 0,20$) e \textbf{Irreversibilidade} ($w_4 = 0,20$). Dois impactos adversos de uma obra rodoviária receberam as seguintes notas:
\begin{center}
\renewcommand{\arraystretch}{1.15}
\small
\begin{tabular}{@{}lcccc@{}}
  \toprule
  \textbf{Impacto ambiental avaliado} & \textbf{Magnitude} ($0,35$) & \textbf{Abrangência} ($0,25$) & \textbf{Duração} ($0,20$) & \textbf{Irreversib.} ($0,20$) \\
  \midrule
  $I_1$: Erosão linear e assoreamento de córrego & $4$ & $3$ & $4$ & $3$ \\
  $I_2$: Elevação temporária de ruído na obra & $2$ & $1$ & $2$ & $1$ \\
  \bottomrule
\end{tabular}
\end{center}
Com base na classificação técnica adotada pelo estudo (Baixa: $1,00 \le I_S < 2,00$; Média: $2,00 \le I_S < 3,00$; e Alta: $3,00 \le I_S \le 4,00$), os valores exatos de $I_S(I_1)$ e $I_S(I_2)$ e suas respectivas classes de significância são:
\begin{alternativas}
  \item $I_S(I_1) = 3,50$ (Alta) e $I_S(I_2) = 1,50$ (Baixa).
  \item $I_S(I_1) = 3,45$ (Alta) e $I_S(I_2) = 1,65$ (Baixa).
  \item $I_S(I_1) = 3,60$ (Alta) e $I_S(I_2) = 1,45$ (Baixa).
  \item $I_S(I_1) = 2,95$ (Média) e $I_S(I_2) = 1,55$ (Baixa).
  \item $I_S(I_1) = 3,55$ (Alta) e $I_S(I_2) = 1,55$ (Baixa).
\end{alternativas}

\vspace{6pt}
\noindent\textbf{Questão 06} {\footnotesize\color{uefsMuted}[Inédita autoral — estilo ENADE / Geografia $\bullet$ Cap. 6]}\par
Ao analisar sistemas ambientais abertos (Bertalanffy, 2010; Drew, 1989) e a gestão do espaço geográfico (Arruda, 2013; Castro; Lemos, 2016), um relatório técnico sobre uma microbacia do Recôncavo Baiano descreveu o seguinte processo: a remoção da vegetação em encostas reduziu a infiltração, elevou o escoamento superficial concentrado, ampliou a erosão e reduziu ainda mais a retenção de umidade do solo, dificultando a regeneração natural e intensificando as perdas de solo nas chuvas seguintes. Esse encadeamento caracteriza um mecanismo de \textbf{retroalimentação} e demanda uma distinção clara entre planejamento e ordenamento territorial. Assinale a alternativa \textbf{correta}.
\begin{alternativas}
  \item Trata-se de retroalimentação \textbf{positiva} (que amplifica o desequilíbrio inicial do sistema); além disso, o planejamento territorial corresponde ao processo contínuo de diagnóstico, definição de diretrizes e metas, enquanto o ordenamento territorial expressa a normatização espacial dos usos e ocupações no território.
  \item Trata-se de retroalimentação \textbf{negativa}, pois o sistema sempre corrige espontaneamente o desmatamento em escala de evento chuvoso; ademais, planejamento e ordenamento territorial são sinônimos estritos sem diferença operacional.
  \item Trata-se de um sistema \textbf{fechado} sem troca de energia ou matéria, no qual a escala espacial de parcela ($1\ \text{m}^2$) é idêntica à escala temporal geológica de evolução de biomas.
  \item Trata-se de retroalimentação \textbf{positiva}, o que significa que o processo é ambientalmente benéfico e indica maior resiliência do geossistema frente às chuvas intensas.
  \item Trata-se de retroalimentação \textbf{negativa}, pois a perda de solo reduz a altitude da vertente, sendo o ordenamento territorial um instrumento exclusivo do setor privado sem vinculação ao poder público.
\end{alternativas}

% ----------------------------------------------------------------------------
\subsection*{Módulo 3 — Bacias, Morfopedologia, Geossistemas e Fragilidade (Capítulos 7 a 9)}

\noindent\textbf{Questão 07} {\footnotesize\color{uefsMuted}[Inédita autoral — estilo Concurso Geógrafo / Recursos Hídricos $\bullet$ Cap. 7]}\par
No diagnóstico hidrogeomorfológico de uma microbacia rural utilizada para abastecimento local, levantaram-se os seguintes parâmetros morfométricos (Christofoletti, 1981; Botelho; Silva, 2004):
\begin{itemize}[leftmargin=1.8em, itemsep=1.5pt, topsep=3pt]
  \item Área da bacia: $A = 25,0\ \text{km}^2$; \quad Perímetro: $P = 22,0\ \text{km}$;
  \item Comprimento total da rede de drenagem: $L_t = 60,0\ \text{km}$;
  \item Altitude máxima: $H_{\max} = 620\ \text{m}$; \quad Altitude mínima no exutório: $H_{\min} = 260\ \text{m}$;
  \item Comprimento axial da bacia: $L_a = 8,0\ \text{km}$.
\end{itemize}
Empregando as expressões clássicas $D_d = L_t / A$, $K_c = 0,28\,P/\sqrt{A}$, $I_c = 12,57\,A/P^2$ e $R_r = (H_{\max} - H_{\min})/L_a$, os valores calculados e a interpretação hidrológica correta sobre a forma dessa microbacia são:
\begin{alternativas}
  \item $D_d = 0,42\ \text{km/km}^2$, $K_c = 1,75$, $I_c = 0,32$ e $R_r = 18,0\ \text{m/km}$, indicando bacia muito alongada e baixa propensão a picos rápidos de cheia.
  \item $D_d = 2,40\ \text{km/km}^2$, $K_c = 1,68$, $I_c = 0,38$ e $R_r = 36,0\ \text{m/km}$, indicando bacia fortemente alongada com escoamento lento e disperso.
  \item $D_d = 2,40\ \text{km/km}^2$, $K_c = 1,23$, $I_c = 0,65$ e $R_r = 45,0\ \text{m/km}$, indicando bacia de formato tendendo a compacto/circular e maior suscetibilidade a picos rápidos de cheia.
  \item $D_d = 2,73\ \text{km/km}^2$, $K_c = 1,23$, $I_c = 0,65$ e $R_r = 77,5\ \text{m/km}$, indicando ausência de risco de cheias rápidas por causa do valor de $K_c$ superior a $1,0$.
  \item $D_d = 2,40\ \text{km/km}^2$, $K_c = 0,81$, $I_c = 0,92$ e $R_r = 45,0\ \text{m/km}$, sendo fisicamente possível obter $K_c < 1,0$ em bacias naturais bem drenadas.
\end{alternativas}

\vspace{6pt}
\noindent\textbf{Questão 08} {\footnotesize\color{uefsMuted}[Inédita autoral — estilo Concurso Analista Ambiental / Geografia Física $\bullet$ Cap. 8]}\par
Em um levantamento morfopedológico de vertente (Casseti, 2007; Guerra, 2018; Ross, 2006), mediu-se um segmento retilíneo de meia encosta que apresenta desnível altimétrico $\Delta H = 90\ \text{m}$ ao longo de uma distância horizontal projetada $L_h = 300\ \text{m}$. Sabendo que $\arctan(0,30) \approx 16,7^\circ$, assinale a alternativa que apresenta corretamente a declividade percentual $S(\%)$, o ângulo de inclinação $\theta$, a classe de relevo (EMBRAPA) e o comportamento morfopedológico esperado nesse segmento de encosta.
\begin{alternativas}
  \item $S = 3,3\%$, $\theta = 30,0^\circ$, relevo suave ondulado, predomínio de solos muito profundos (Latossolos) e infiltração máxima sem risco erosivo.
  \item $S = 30,0\%$, $\theta = 30,0^\circ$, relevo montanhoso, predomínio de Gleissolos hidromórficos em função da estagnação de água na meia encosta.
  \item $S = 30,0\%$, $\theta = 45,0^\circ$, relevo escarpado, no qual a pedogênese supera amplamente a morfogênese independentemente da cobertura vegetal.
  \item $S = 16,7\%$, $\theta = 30,0^\circ$, relevo ondulado, com baixa velocidade de escoamento superficial e deposição predominante de aluviões arenosos.
  \item $S = 30,0\%$, $\theta \approx 16,7^\circ$, relevo forte ondulado ($20\%$ a $45\%$), com alta energia de escoamento superficial, tendência a solos menos espessos (Neossolos/Cambissolos) e elevado potencial erosivo quando desprotegido.
\end{alternativas}

\vspace{6pt}
\noindent\textbf{Questão 09} {\footnotesize\color{uefsMuted}[Inédita autoral — estilo ENADE / Concurso Geógrafo $\bullet$ Cap. 9]}\par
Na metodologia de Fragilidade Ambiental de Jurandyr Ross (1994), fundamentada nos princípios da Ecodinâmica de Jean Tricart (1977), a \textbf{Fragilidade Potencial (FP)} expressa a vulnerabilidade intrínseca do meio físico (relevo/declividade, solos e clima), enquanto a \textbf{Fragilidade Emergente (FE)} integra o grau de proteção conferido pelo uso e cobertura da terra. Considere uma vertente no Semiárido com as seguintes classes de entrada (escala de $1$ = Muito Fraca a $5$ = Muito Forte):
\begin{itemize}[leftmargin=1.8em, itemsep=1.5pt, topsep=3pt]
  \item Declividade entre $20\%$ e $30\%$ $\to$ \textbf{Classe 4} (Forte);
  \item Solo raso e pedregoso (Neossolo Litólico) $\to$ \textbf{Classe 5} (Muito Forte);
  \item Regime pluviométrico concentrado em poucos meses $\to$ \textbf{Classe 3} (Média).
\end{itemize}
Adotando a média aritmética das três variáveis físicas para estimar a Fragilidade Potencial ($FP$) e cruzando o resultado com a matriz de Fragilidade Emergente de Ross (em que $FP = 4$ sob \textbf{Mata densa} [proteção classe 1] resulta em $FE = 2$, enquanto $FP = 4$ sob \textbf{Solo exposto} [proteção classe 5] resulta em $FE = 5$), assinale a afirmação \textbf{correta}.
\begin{alternativas}
  \item A Fragilidade Potencial é igual a $3$ (Média), razão pela qual a retirada total da cobertura vegetal não altera o balanço entre pedogênese e morfogênese.
  \item A Fragilidade Potencial é igual a $4$ (Forte); sob solo exposto, a Fragilidade Emergente atinge a classe $5$ (Muito Forte), caracterizando um meio fortemente instável na concepção de Tricart, com nítido predomínio da morfogênese sobre a pedogênese.
  \item A Fragilidade Emergente é uma propriedade exclusivamente natural que independe do manejo agrícola, sendo igual a $4$ tanto sob mata densa quanto sob solo exposto.
  \item Na Ecodinâmica de Tricart (1977), o cenário de solo exposto com $FE = 5$ corresponde a um meio estável, pois a erosão rápida estabiliza imediatamente o perfil pedológico.
  \item A conservação de mata densa sobre essa vertente eleva a Fragilidade Potencial para classe $5$, pois o peso das árvores desestabiliza os Neossolos Litólicos.
\end{alternativas}

% ----------------------------------------------------------------------------
\subsection*{Módulo 4 — Indicadores, Qualidade do Solo e Degradação (Capítulos 10 a 12)}

\noindent\textbf{Questão 10} {\footnotesize\color{uefsMuted}[Inédita autoral — estilo Concurso Analista Ambiental / Semiárido $\bullet$ Cap. 10]}\par
Durante o trabalho de campo no eixo energético do Semiárido baiano (região de Canudos e complexos eólicos na Caatinga), os estudantes organizaram um painel de monitoramento ambiental segundo o modelo \textbf{PEIR} (OECD, 2003; IBGE, 2018; Fidalgo, 2003), listando quatro indicadores:
\begin{enumerate}[label=\arabic*., leftmargin=2.0em, itemsep=1.5pt, topsep=3pt]
  \item Extensão (em $\text{km}$) de novos acessos abertos e área de Caatinga suprimida para platôs de aerogeradores;
  \item Percentual de cobertura vegetal remanescente e densidade de solo exposto na microbacia;
  \item Volume de sedimentos erodidos na base das torres e taxa de mortalidade de avifauna por colisão;
  \item Percentual de áreas degradadas efetivamente recuperadas via PRAD, drenagem e bioengenharia.
\end{enumerate}
A correspondência \textbf{correta} dos indicadores 1, 2, 3 e 4 às categorias do modelo PEIR é, respectivamente:
\begin{alternativas}
  \item Estado -- Pressão -- Resposta -- Impacto.
  \item Impacto -- Estado -- Pressão -- Resposta.
  \item Pressão -- Impacto -- Estado -- Resposta.
  \item Pressão -- Estado -- Impacto -- Resposta.
  \item Resposta -- Pressão -- Estado -- Impacto.
\end{alternativas}

\vspace{6pt}
\noindent\textbf{Questão 11} {\footnotesize\color{uefsMuted}[Inédita autoral — estilo Concurso Analista de Solos / IBAMA $\bullet$ Cap. 11]}\par
Para avaliar a Qualidade do Solo (Doran; Parkin, 1994; Karlen et al., 1997; Vezzani; Mielniczuk, 2009) em uma pastagem degradada com pisoteio intensivo, coletou-se uma amostra indeformada com anel volumétrico de volume $V = 100,0\ \text{cm}^3$, cuja massa de solo seco em estufa a $105\,^\circ\text{C}$ foi $m_s = 165,0\ \text{g}$. Em análise química complementar, determinou-se o teor de Carbono Orgânico total $CO = 10,0\ \text{g/kg}$. Adotando a densidade de partículas padrão $D_p = 2,65\ \text{g/cm}^3$, as relações $D_s = m_s / V$, $P_t(\%) = (1 - D_s / D_p)\times 100$ e o fator de conversão de Van Bemmelen ($MO = 1,724 \times CO$), os valores obtidos e o diagnóstico físico do solo em comparação a uma área de referência de Caatinga conservada ($D_{s,\text{ref}} = 1,15\ \text{g/cm}^3$; $P_{t,\text{ref}} = 56,60\%$) são:
\begin{alternativas}
  \item $D_s = 1,65\ \text{g/cm}^3$, $P_t = 62,26\%$ e $MO = 10,00\ \text{g/kg}$, indicando excelente aeração e ausência de compactação.
  \item $D_s = 1,65\ \text{g/cm}^3$, $P_t = 37,74\%$ e $MO = 17,24\ \text{g/kg}$, evidenciando forte compactação e redução de $18,87$ pontos percentuais na porosidade total em relação à área de referência.
  \item $D_s = 0,61\ \text{g/cm}^3$, $P_t = 37,74\%$ e $MO = 5,80\ \text{g/kg}$, indicando solo extremamente poroso e rico em matéria orgânica.
  \item $D_s = 1,65\ \text{g/cm}^3$, $P_t = 43,40\%$ e $MO = 17,24\ \text{g/kg}$, revelando ganho de macroporosidade provocado pelo pisoteio animal.
  \item $D_s = 2,65\ \text{g/cm}^3$, $P_t = 37,74\%$ e $MO = 26,50\ \text{g/kg}$, demonstrando que a densidade do solo independe do arranjo estrutural dos poros.
\end{alternativas}

\vspace{6pt}
\noindent\textbf{Questão 12} {\footnotesize\color{uefsMuted}[Inédita autoral — estilo Concurso Engenharia / Geografia Ambiental $\bullet$ Cap. 12]}\par
Em uma parcela agrícola inclinada no interior baiano, estimou-se a perda média anual de solo pela Equação Universal de Perda de Solo ($A = R \cdot K \cdot L \cdot S \cdot C \cdot P$) sob dois cenários de manejo, mantendo-se constantes os fatores naturais: erosividade da chuva $R = 5.000\ \text{MJ}\cdot\text{mm}\cdot\text{ha}^{-1}\cdot\text{h}^{-1}\cdot\text{ano}^{-1}$, erodibilidade do solo $K = 0,030\ \text{t}\cdot\text{h}\cdot\text{MJ}^{-1}\cdot\text{mm}^{-1}$ e fator topográfico combinado $LS = 2,0$.
\begin{itemize}[leftmargin=1.8em, itemsep=1.5pt, topsep=3pt]
  \item \textbf{Cenário 1 (atual — cultivo convencional morro abaixo):} $C_1 = 0,40$ e $P_1 = 1,00$;
  \item \textbf{Cenário 2 (proposto — plantio direto com cobertura morta e plantio em nível):} $C_2 = 0,05$ e $P_2 = 0,50$.
\end{itemize}
Considerando um limite de tolerância de perda de solo $T = 12,0\ \text{t}\cdot\text{ha}^{-1}\cdot\text{ano}^{-1}$ e as limitações conceituais da USLE (Guerra; Cunha, 1996), é \textbf{correto} afirmar que:
\begin{alternativas}
  \item a perda de solo cai de $120,0\ \text{t}\cdot\text{ha}^{-1}\cdot\text{ano}^{-1}$ para $7,5\ \text{t}\cdot\text{ha}^{-1}\cdot\text{ano}^{-1}$ (redução de $93,75\%$), trazendo a parcela para abaixo da tolerância $T$; contudo, a USLE estima erosão laminar e em sulcos, não modelando diretamente o colapso de voçorocas por \textit{piping}.
  \item a perda de solo cai de $300,0\ \text{t}\cdot\text{ha}^{-1}\cdot\text{ano}^{-1}$ para $15,0\ \text{t}\cdot\text{ha}^{-1}\cdot\text{ano}^{-1}$ (redução de $95,0\%$), permanecendo acima da tolerância $T$.
  \item a perda de solo cai de $120,0\ \text{t}\cdot\text{ha}^{-1}\cdot\text{ano}^{-1}$ para $15,0\ \text{t}\cdot\text{ha}^{-1}\cdot\text{ano}^{-1}$ (redução de $87,5\%$), pois o fator $P$ não multiplica o produto dos demais fatores.
  \item a perda de solo cai de $60,0\ \text{t}\cdot\text{ha}^{-1}\cdot\text{ano}^{-1}$ para $7,5\ \text{t}\cdot\text{ha}^{-1}\cdot\text{ano}^{-1}$, e a USLE calcula exatamente o volume de sedimentos produzido por solapamento do lençol freático em voçorocas.
  \item a perda de solo aumenta com a redução dos fatores $C$ e $P$, porque o plantio em nível concentra o escoamento superficial ao longo da linha de maior declive.
\end{alternativas}

% ----------------------------------------------------------------------------
\subsection*{Módulo 5 — Bioengenharia, Legislação, PRAD e Diagnóstico Final (Capítulos 13 a 15)}

\noindent\textbf{Questão 13} {\footnotesize\color{uefsMuted}[Inédita autoral — estilo Concurso Analista Ambiental / Recuperação de Áreas $\bullet$ Cap. 13]}\par
Para estabilizar uma ravina com comprimento longitudinal $L = 60,0\ \text{m}$, largura média $b = 2,0\ \text{m}$ e declividade longitudinal do leito $S = 10\%$ ($0,10\ \text{m/m}$), um projeto de bioengenharia de água e solo (Gray; Sotir, 1996; Schiechtl; Stern, 1996; Durlo; Sutili, 2005) prevê a instalação de paliçadas permeáveis de bambu e madeira tratada com altura útil do vertedouro $h = 0,60\ \text{m}$, combinadas a estacas vivas. Utilizando o critério geométrico básico de espaçamento horizontal $E = h / S$ (de modo que o topo da cunha de sedimento retido alcance o pé da paliçada imediatamente a montante) e aproximando cada cunha retida por um prisma triangular de volume $V_{\text{cunha}} = \tfrac{1}{2}\,b\,h\,E$, o espaçamento $E$, o número de trechos/estruturas ao longo dos $60\ \text{m}$ ($n = L / E$) e a capacidade total de retenção de sedimentos ($V_{\text{total}} = n \cdot V_{\text{cunha}}$) são, respectivamente:
\begin{alternativas}
  \item $E = 0,60\ \text{m}$, $n = 100$ paliçadas e $V_{\text{total}} = 36,0\ \text{m}^3$.
  \item $E = 3,00\ \text{m}$, $n = 20$ paliçadas e $V_{\text{total}} = 18,0\ \text{m}^3$.
  \item $E = 6,00\ \text{m}$, $n = 10$ paliçadas e $V_{\text{total}} = 72,0\ \text{m}^3$.
  \item $E = 12,00\ \text{m}$, $n = 5$ paliçadas e $V_{\text{total}} = 36,0\ \text{m}^3$.
  \item $E = 6,00\ \text{m}$, $n = 10$ paliçadas e $V_{\text{total}} = 36,0\ \text{m}^3$.
\end{alternativas}

\vspace{6pt}
\noindent\textbf{Questão 14} {\footnotesize\color{uefsMuted}[Inédita autoral — estilo Concurso IBAMA / Direito e Gestão Ambiental $\bullet$ Cap. 14]}\par
Um imóvel rural de $150,0\ \text{ha}$, localizado no bioma \textbf{Caatinga} (fora da Amazônia Legal), é integralmente atravessado por um curso d'água natural perene cuja calha regular mede $35\ \text{m}$ de largura, ao longo de uma extensão de $1.200\ \text{m}$ (com ambas as margens situadas dentro da propriedade). À luz do Código Florestal Brasileiro (Lei nº 12.651/2012, Arts. 4º e 12) e sem considerar o cômputo excepcional da APP na Reserva Legal previsto no Art. 15, a largura mínima da Área de Preservação Permanente (APP) em cada margem, a área total de APP marginal nesse trecho e a área mínima de Reserva Legal (RL) exigida para o imóvel correspondem, respectivamente, a:
\begin{alternativas}
  \item $30\ \text{m}$ em cada margem; $7,2\ \text{ha}$ de APP marginal; e $30,0\ \text{ha}$ de Reserva Legal ($20\%$).
  \item $50\ \text{m}$ em cada margem; $6,0\ \text{ha}$ de APP marginal; e $52,5\ \text{ha}$ de Reserva Legal ($35\%$).
  \item $50\ \text{m}$ em cada margem; $12,0\ \text{ha}$ de APP marginal; e $30,0\ \text{ha}$ de Reserva Legal ($20\%$).
  \item $100\ \text{m}$ em cada margem; $24,0\ \text{ha}$ de APP marginal; e $30,0\ \text{ha}$ de Reserva Legal ($20\%$).
  \item $50\ \text{m}$ em cada margem; $12,0\ \text{ha}$ de APP marginal; e $120,0\ \text{ha}$ de Reserva Legal ($80\%$).
\end{alternativas}

\vspace{6pt}
\noindent\textbf{Questão 15} {\footnotesize\color{uefsMuted}[Inédita autoral — estilo ENADE / Trabalho Prático Profissional $\bullet$ Cap. 15]}\par
Na elaboração de um Diagnóstico Ambiental Aplicado e de um Plano de Recuperação de Áreas Degradadas (PRAD) para uma microbacia degradada (Vital; Carvalho; Pinto, 2020; IBGE, 2018; Sánchez, 2013), a coerência metodológica entre diagnóstico, indicadores e intervenção é condição decisiva para aprovação em órgãos licenciadores. Considere uma equipe que diagnosticou voçorocamento em cabeceira de drenagem, solo compactado ($D_s = 1,68\ \text{g/cm}^3$) e supressão de APP ripária. A proposta técnica \textbf{mais consistente} do ponto de vista da Análise Ambiental integrada é:
\begin{alternativas}
  \item plantar diretamente mudas arbóreas de crescimento lento dentro da voçoroca ativa sem desvio das águas pluviais a montante, dispensando monitoramento pós-plantio.
  \item apresentar apenas um levantamento fotográfico qualitativo sem delimitação de bacia, sem indicadores mensuráveis e sem cronograma físico de execução.
  \item substituir integralmente o solo natural por concreto armado em toda a microbacia, desconsiderando a recarga hídrica e a legislação de APPs.
  \item articular recorte espacial da microbacia, controle prévio do escoamento a montante (terraços/canais divergentes), obras de bioengenharia (paliçadas permeáveis, retaludamento suave, biomanta e revegetação com espécies nativas adaptadas) e matriz de indicadores monitoráveis de solo, sedimentos e sobrevivência vegetal.
  \item limitar o PRAD à compensação financeira em outro município, mantendo aberta a erosão ativa na cabeceira de drenagem sem medidas de contenção local.
\end{alternativas}

% ============================================================================
\clearpage
\section*{Parte II — Questões Dissertativas Integradoras (Padrão Prova UEFS)}
\addcontentsline{toc}{section}{Parte II — Questões Dissertativas Integradoras (D1–D5)}
% ============================================================================

\noindent Responda às cinco questões dissertativas a seguir com rigor conceitual, clareza argumentativa e articulação explícita entre os fundamentos teóricos e os instrumentos técnicos trabalhados na disciplina CHF662 (0,8 ponto cada; total: 4,0 pontos).

\vspace{6pt}
\noindent\textbf{Dissertativa 1 (Módulo 1 — Epistemologia, Sociedade--Natureza e Conceitos-Chave):}\par
Explique por que a Geografia Socioambiental (Mendonça, 2001) e a crítica epistemológica de Leff (2001) e Porto-Gonçalves (2006) rejeitam a visão naturalista estrita dos problemas ambientais. Em seguida, diferencie com precisão \textbf{perigo}, \textbf{vulnerabilidade}, \textbf{risco} e \textbf{resiliência ambiental}, ilustrando com uma situação concreta do Semiárido ou de áreas urbanas baianas.
\espacoResolucao[3.2cm]

\noindent\textbf{Dissertativa 2 (Módulo 2 — AIA, Escalas e Planejamento Ambiental):}\par
Descreva o encadeamento lógico das etapas da Avaliação de Impacto Ambiental (Sánchez, 2013), destacando a diferença funcional entre \textbf{diagnóstico ambiental} (linha de base) e \textbf{prognóstico} (comparação de cenários com e sem projeto). Justifique por que o descompasso entre escalas espaciais e temporais pode invalidar um Estudo de Impacto Ambiental e comprometer o ordenamento territorial.
\espacoResolucao[3.2cm]

\noindent\textbf{Dissertativa 3 (Módulo 3 — Bacias Hidrográficas, Morfopedologia e Fragilidade):}\par
Analise por que a \textbf{bacia hidrográfica} constitui a unidade espacial mais adequada para a Análise Ambiental integrada (Botelho; Silva, 2004). Articule em sua resposta: (i) a variação morfopedológica ao longo de uma vertente (topo, meia encosta, sopé e fundo de vale); (ii) o balanço morfogênese versus pedogênese de Tricart (1977); e (iii) a passagem da Fragilidade Potencial para a Fragilidade Emergente no método de Ross (1994).
\espacoResolucao[3.2cm]

\noindent\textbf{Dissertativa 4 (Módulo 4 — Indicadores PEIR, Qualidade do Solo e Erosão Linear):}\par
Discuta a importância de integrar indicadores físicos, químicos e biológicos na avaliação da Qualidade do Solo (Doran; Parkin, 1994; Vezzani; Mielniczuk, 2009), explicando como a perda de matéria orgânica e o aumento da densidade do solo desencadeiam a sequência evolutiva da erosão hídrica (do salpicamento à voçoroca). Indique como o modelo PEIR pode organizar o monitoramento de uma área degradada na Caatinga.
\espacoResolucao[3.2cm]

\noindent\textbf{Dissertativa 5 (Módulo 5 — Bioengenharia, Legislação Ambiental e PRAD):}\par
Compare os efeitos \textbf{hidrológicos} e \textbf{mecânicos} da vegetação na estabilização de taludes e margens fluviais (Gray; Sotir, 1996; Durlo; Sutili, 2005), referindo-se explicitamente à coesão radicular ($c_r$) no critério de Mohr--Coulomb. Em seguida, relacione a aplicação de técnicas de bioengenharia às exigências legais de recuperação de APPs no Código Florestal (Lei nº 12.651/2012) e às fases do licenciamento ambiental trifásico (LP, LI e LO).
\espacoResolucao[3.2cm]

% ============================================================================
\clearpage
\section*{Folha de Respostas Imprimível — Parte I (Questões Q01 a Q15)}
\addcontentsline{toc}{section}{Folha de Respostas Imprimível (Q01–Q15)}
% ============================================================================

\noindent Preencha totalmente a letra escolhida para cada questão antes de conferir o Gabarito Comentado na página seguinte.

\vspace{6pt}
\begin{center}
\renewcommand{\arraystretch}{1.45}
\small
\begin{tabularx}{\textwidth}{|c|c|*{5}{>{\centering\arraybackslash}X|}c|}
  \hline
  \rowcolor{uefsNavy}
  \textcolor{white}{\bfseries Questão} & \textcolor{white}{\bfseries Módulo / Cap.} &
  \textcolor{white}{\bfseries A} & \textcolor{white}{\bfseries B} & \textcolor{white}{\bfseries C} & \textcolor{white}{\bfseries D} & \textcolor{white}{\bfseries E} &
  \textcolor{white}{\bfseries Acerto?} \\
  \hline
  \textbf{Q01} & Mód. 1 (Caps. 1--2) & ( \quad ) & ( \quad ) & ( \quad ) & ( \quad ) & ( \quad ) & [ \quad ] \\ \hline
  \rowcolor{bgGray}
  \textbf{Q02} & Mód. 1 (Cap. 3)     & ( \quad ) & ( \quad ) & ( \quad ) & ( \quad ) & ( \quad ) & [ \quad ] \\ \hline
  \textbf{Q03} & Mód. 1 (Cap. 4)     & ( \quad ) & ( \quad ) & ( \quad ) & ( \quad ) & ( \quad ) & [ \quad ] \\ \hline
  \rowcolor{bgGray}
  \textbf{Q04} & Mód. 2 (Cap. 5)     & ( \quad ) & ( \quad ) & ( \quad ) & ( \quad ) & ( \quad ) & [ \quad ] \\ \hline
  \textbf{Q05} & Mód. 2 (Cap. 5)     & ( \quad ) & ( \quad ) & ( \quad ) & ( \quad ) & ( \quad ) & [ \quad ] \\ \hline
  \rowcolor{bgGray}
  \textbf{Q06} & Mód. 2 (Cap. 6)     & ( \quad ) & ( \quad ) & ( \quad ) & ( \quad ) & ( \quad ) & [ \quad ] \\ \hline
  \textbf{Q07} & Mód. 3 (Cap. 7)     & ( \quad ) & ( \quad ) & ( \quad ) & ( \quad ) & ( \quad ) & [ \quad ] \\ \hline
  \rowcolor{bgGray}
  \textbf{Q08} & Mód. 3 (Cap. 8)     & ( \quad ) & ( \quad ) & ( \quad ) & ( \quad ) & ( \quad ) & [ \quad ] \\ \hline
  \textbf{Q09} & Mód. 3 (Cap. 9)     & ( \quad ) & ( \quad ) & ( \quad ) & ( \quad ) & ( \quad ) & [ \quad ] \\ \hline
  \rowcolor{bgGray}
  \textbf{Q10} & Mód. 4 (Cap. 10)    & ( \quad ) & ( \quad ) & ( \quad ) & ( \quad ) & ( \quad ) & [ \quad ] \\ \hline
  \textbf{Q11} & Mód. 4 (Cap. 11)    & ( \quad ) & ( \quad ) & ( \quad ) & ( \quad ) & ( \quad ) & [ \quad ] \\ \hline
  \rowcolor{bgGray}
  \textbf{Q12} & Mód. 4 (Cap. 12)    & ( \quad ) & ( \quad ) & ( \quad ) & ( \quad ) & ( \quad ) & [ \quad ] \\ \hline
  \textbf{Q13} & Mód. 5 (Cap. 13)    & ( \quad ) & ( \quad ) & ( \quad ) & ( \quad ) & ( \quad ) & [ \quad ] \\ \hline
  \rowcolor{bgGray}
  \textbf{Q14} & Mód. 5 (Cap. 14)    & ( \quad ) & ( \quad ) & ( \quad ) & ( \quad ) & ( \quad ) & [ \quad ] \\ \hline
  \textbf{Q15} & Mód. 5 (Cap. 15)    & ( \quad ) & ( \quad ) & ( \quad ) & ( \quad ) & ( \quad ) & [ \quad ] \\ \hline
\end{tabularx}
\end{center}

\vspace{10pt}
\begin{tcolorbox}[colback=bgGold, colframe=uefsGold, arc=3pt, boxrule=0.7pt, title={\bfseries Registro Rápido de Pontuação Final}, coltitle=white, colbacktitle=uefsGold]
\small
\begin{tabularx}{\textwidth}{@{}X X X@{}}
  \textbf{Acertos Parte I:} \underline{\hspace{1.2cm}} / 15 &
  \textbf{Nota Parte I} ($\times 0,4$): \underline{\hspace{1.2cm}} / 6,0 &
  \textbf{Nota Parte II:} \underline{\hspace{1.2cm}} / 4,0 \\[6pt]
  \multicolumn{3}{@{}l}{\textbf{Nota Geral Estimada no Simulado (Parte I + Parte II):} \underline{\hspace{2.0cm}} / 10,0}
\end{tabularx}
\end{tcolorbox}

% ============================================================================
\clearpage
\section*{Gabarito Integralmente Comentado — Parte I (Q01 a Q15)}
\addcontentsline{toc}{section}{Gabarito Comentado, Espelho de Correção e Diagnóstico Final}
% ============================================================================

\begin{center}
\renewcommand{\arraystretch}{1.25}
\small
\begin{tabular}{|*{15}{c|}}
  \hline
  \rowcolor{uefsNavy}
  \textcolor{white}{\bfseries 01} & \textcolor{white}{\bfseries 02} & \textcolor{white}{\bfseries 03} &
  \textcolor{white}{\bfseries 04} & \textcolor{white}{\bfseries 05} & \textcolor{white}{\bfseries 06} &
  \textcolor{white}{\bfseries 07} & \textcolor{white}{\bfseries 08} & \textcolor{white}{\bfseries 09} &
  \textcolor{white}{\bfseries 10} & \textcolor{white}{\bfseries 11} & \textcolor{white}{\bfseries 12} &
  \textcolor{white}{\bfseries 13} & \textcolor{white}{\bfseries 14} & \textcolor{white}{\bfseries 15} \\
  \hline
  \textbf{C} & \textbf{A} & \textbf{D} & \textbf{B} & \textbf{E} &
  \textbf{A} & \textbf{C} & \textbf{E} & \textbf{B} & \textbf{D} &
  \textbf{B} & \textbf{A} & \textbf{E} & \textbf{C} & \textbf{D} \\
  \hline
\end{tabular}
\end{center}

\begin{boxGabarito}[Resolução Passo a Passo e Análise dos Distratores (Q01 a Q15)]
\small
\begin{enumerate}[leftmargin=1.8em, itemsep=5pt]
  \item \textbf{Gabarito: (C).} Na epistemologia ambiental (Leff, 2001) e na crítica à modernidade de Porto-Gonçalves (2006), a crise ambiental é uma crise de civilização e de racionalidade, não um defeito meramente técnico ou climático. Atribuir conflitos hídricos exclusivamente à seca naturaliza desigualdades de acesso e oculta interesses econômicos e territoriais. \textit{Distratores:} (A) e (D) reforçam o reducionismo fisicalista e a dicotomia geográfica; (B) reduz a questão ao plano jurídico e defende solução única tecnocrática; (E) distorce o argumento imputando controle climático às comunidades.

  \item \textbf{Gabarito: (A).} No sistema GTP (Bertrand, 2004), \textit{Geossistema} trata da dinâmica natural e das fontes físico-ecológicas; \textit{Território} aborda a apropriação social, econômica e político-institucional; e \textit{Paisagem} expressa a dimensão cultural, simbólica e perceptiva do espaço. \textit{Distratores:} (B) e (C) restringem ou invertem os três polos; (D) troca o conteúdo de Geossistema e Paisagem; (E) afirma erroneamente que o geossistema elimina território e região.

  \item \textbf{Gabarito: (D).} Aplicando $R = P \times V$:
  \begin{itemize}[leftmargin=1.4em, itemsep=1pt, topsep=2pt]
    \item Setor A (inicial): $R_A = 4 \times 3 = 12$ (\textbf{Crítico}, intervalo $12$ a $16$);
    \item Setor B: $R_B = 3 \times 2 = 6$ (\textbf{Médio}, intervalo $4$ a $6$);
    \item Setor C: $R_C = 2 \times 4 = 8$ (\textbf{Alto}, intervalo $8$ a $9$);
    \item Setor A (após redução de $V$ para $1$): $R_{A,\text{novo}} = 4 \times 1 = 4$ (\textbf{Médio});
    \item Redução relativa: $\frac{12 - 4}{12} \times 100 = \frac{8}{12} \times 100 \approx 66,7\%$.
  \end{itemize}
  \textit{Distratores:} (A) erra ao atribuir peso quadrático a $V$ e ao inverter a ordem de A e C; (B) inverte atividade/aspecto com impacto (corte de talude é atividade, desestabilização/exposição é aspecto, deslizamento sobre moradias é impacto); (C) erra o produto $4 \times 1 = 4$; (E) confunde aspecto (exposição do solo) e impacto (assoreamento) com atividades primárias.

  \item \textbf{Gabarito: (B).} A hierarquia de mitigação prioriza: 1º \textbf{Evitar} na origem (I — mudar traçado para desviar da vereda e da mata ciliar); 2º \textbf{Minimizar} a magnitude (III — cabos por drones e controle de drenagem); 3º \textbf{Recuperar/Reabilitar} áreas afetadas (IV — PRAD e bioengenharia); e 4º \textbf{Compensar} somente o impacto residual inevitável (II — apoio à Unidade de Conservação). Portanto: \textbf{I -- III -- IV -- II}. \textit{Distratores:} (A), (C), (D) e (E) subvertem a ordem preventiva ao colocar minimização, recuperação ou compensação antes de evitar o dano.

  \item \textbf{Gabarito: (E).} Calculando $I_S = 0,35\,n_1 + 0,25\,n_2 + 0,20\,n_3 + 0,20\,n_4$:
  \begin{itemize}[leftmargin=1.4em, itemsep=1pt, topsep=2pt]
    \item $I_S(I_1) = (0,35 \times 4) + (0,25 \times 3) + (0,20 \times 4) + (0,20 \times 3) = 1,40 + 0,75 + 0,80 + 0,60 = 3,55$ (\textbf{Alta}, pois $3,00 \le 3,55 \le 4,00$);
    \item $I_S(I_2) = (0,35 \times 2) + (0,25 \times 1) + (0,20 \times 2) + (0,20 \times 1) = 0,70 + 0,25 + 0,40 + 0,20 = 1,55$ (\textbf{Baixa}, pois $1,00 \le 1,55 < 2,00$).
  \end{itemize}
  \textit{Distratores:} (A) corresponde à média aritmética simples sem ponderação ($14/4 = 3,50$ e $6/4 = 1,50$); (B), (C) e (D) resultam de trocas nos pesos atribuídos a Magnitude e Abrangência.

  \item \textbf{Gabarito: (A).} Quando a saída do sistema (erosão e perda de infiltração) reforça a causa inicial e intensifica a degradação, tem-se \textbf{retroalimentação positiva} (autorreforço, não juízo de valor ``bom''). Na distinção de Arruda (2013) e Castro e Lemos (2016), planejamento territorial é o processo técnico-político de diagnóstico e definição de diretrizes, enquanto ordenamento territorial corresponde à materialização normativa da distribuição dos usos no território. \textit{Distratores:} (B) e (E) confundem retroalimentação positiva com negativa e distorcem o papel do Estado; (C) erra ao chamar a bacia de sistema fechado e igualar escalas incompatíveis; (D) confunde ``positivo'' (amplificação matemática) com ``benéfico''.

  \item \textbf{Gabarito: (C).} Calculando cada índice passo a passo:
  \begin{itemize}[leftmargin=1.4em, itemsep=1pt, topsep=2pt]
    \item Densidade de drenagem: $D_d = \frac{L_t}{A} = \frac{60,0}{25,0} = 2,40\ \text{km/km}^2$;
    \item Coeficiente de compacidade: $K_c = 0,28 \cdot \frac{P}{\sqrt{A}} = 0,28 \cdot \frac{22,0}{\sqrt{25,0}} = 0,28 \cdot \frac{22,0}{5,0} = 0,28 \times 4,40 = 1,232 \approx 1,23$;
    \item Índice de circularidade: $I_c = 12,57 \cdot \frac{A}{P^2} = 12,57 \cdot \frac{25,0}{22,0^2} = \frac{314,25}{484,0} \approx 0,649 \approx 0,65$;
    \item Relação de relevo: $R_r = \frac{H_{\max} - H_{\min}}{L_a} = \frac{620 - 260}{8,0} = \frac{360}{8,0} = 45,0\ \text{m/km}$.
  \end{itemize}
  Como $K_c = 1,23$ aproxima-se da unidade ($K_c = 1,0$ seria um círculo perfeito) e $I_c = 0,65 > 0,51$, a bacia tende ao formato compacto/circular, favorecendo concentração simultânea do deflúvio no exutório e maior propensão a cheias rápidas. \textit{Distratores:} (A) inverte $A/L_t$ para $D_d$; (B) erra $K_c$, $I_c$ e $R_r$; (D) usa $H_{\max}/L_a$ em vez de $\Delta H/L_a$ e erra a interpretação de $K_c$; (E) viola o teorema isoperimétrico, pois nenhuma figura plana tem $K_c < 1,0$.

  \item \textbf{Gabarito: (E).} A declividade percentual é a razão entre o desnível vertical e a distância horizontal multiplicada por $100$:
  \[
    S(\%) = \frac{\Delta H}{L_h}\times 100 = \frac{90}{300}\times 100 = 30,0\%, \qquad
    \theta = \arctan\!\left(\frac{90}{300}\right) = \arctan(0,30) \approx 16,7^\circ.
  \]
  Na classificação EMBRAPA, $30,0\%$ situa-se na faixa de relevo \textbf{forte ondulado} ($20\%$ a $45\%$), típico de meias encostas íngremes com escoamento veloz, solos rasos/jovens (Neossolos e Cambissolos) e elevado risco erosivo sem cobertura vegetal. \textit{Distratores:} (A) inverte a razão ($\Delta H/L_h$) e confunde $\%$ com grau; (B) comete o erro clássico de igualar $30\%$ a $30^\circ$ e associa Gleissolos à meia encosta; (C) iguala $30\%$ a $45^\circ$ (que corresponde a $100\%$); (D) troca os valores de $S(\%)$ e $\theta$.

  \item \textbf{Gabarito: (B).} Calculando a média das três variáveis físicas:
  \[
    FP = \frac{4 + 5 + 3}{3} = \frac{12}{3} = 4 \quad (\text{Classe 4 — Forte}).
  \]
  Cruzando $FP = 4$ com \textbf{Solo exposto} (proteção mínima, classe 5) na matriz de Ross (1994), obtém-se $FE = 5$ (\textbf{Muito Forte}), correspondendo na Ecodinâmica de Tricart (1977) a um \textbf{meio fortemente instável}, no qual os processos morfogenéticos (erosão e remoção de material) superam a pedogênese. \textit{Distratores:} (A) erra a média aritmética; (C) confunde Fragilidade Potencial (natural) com Emergente (que inclui o uso da terra); (D) inverte meio estável com fortemente instável; (E) desconhece o papel protetor da cobertura florestal em encostas.

  \item \textbf{Gabarito: (D).} No modelo PEIR (OECD, 2003; IBGE, 2018): (1) abertura de acessos e supressão vegetal são ações antrópicas geradoras de estresse $\to$ \textbf{Pressão (P)}; (2) percentual de cobertura remanescente e solo exposto retratam a condição atual do ambiente $\to$ \textbf{Estado (E)}; (3) volume erodido e mortalidade da avifauna expressam consequências sobre ecossistemas e infraestrutura $\to$ \textbf{Impacto (I)}; e (4) áreas recuperadas por PRAD e bioengenharia representam intervenções corretivas da sociedade $\to$ \textbf{Resposta (R)}. \textit{Distratores:} (A), (B), (C) e (E) permutam incorretamente as categorias causais do ciclo PEIR.

  \item \textbf{Gabarito: (B).} Efetuando cada cálculo passo a passo:
  \begin{itemize}[leftmargin=1.4em, itemsep=1pt, topsep=2pt]
    \item Densidade do solo: $D_s = \frac{m_s}{V} = \frac{165,0\ \text{g}}{100,0\ \text{cm}^3} = 1,65\ \text{g/cm}^3$;
    \item Porosidade total: $P_t = \left(1 - \frac{1,65}{2,65}\right)\times 100 = (1 - 0,62264)\times 100 \approx 37,74\%$;
    \item Matéria orgânica: $MO = 1,724 \times CO = 1,724 \times 10,0 = 17,24\ \text{g/kg}$;
    \item Redução absoluta de porosidade em relação à referência ($P_{t,\text{ref}} = 56,60\%$): $56,60 - 37,74 = 18,86 \approx 18,87$ pontos percentuais (diferença exata usando valores não arredondados: $56,6038 - 37,7358 = 18,8679\ \text{p.p.}$).
  \end{itemize}
  \textit{Distratores:} (A) usa $D_s/D_p = 62,26\%$ em vez de $1 - D_s/D_p$ e omite o fator $1,724$; (C) inverte $V/m_s$ e divide $CO$ por $1,724$; (D) calcula $1 - 1,50/2,65$ e inverte o efeito do pisoteio; (E) confunde $D_s$ com $D_p$.

  \item \textbf{Gabarito: (A).} Primeiro obtém-se o potencial natural de erosão ($R \cdot K \cdot LS$):
  \[
    R \cdot K \cdot LS = 5.000 \times 0,030 \times 2,0 = 300,0\ \text{t}\cdot\text{ha}^{-1}\cdot\text{ano}^{-1}.
  \]
  Aplicando os fatores de uso e manejo nos dois cenários:
  \begin{itemize}[leftmargin=1.4em, itemsep=1pt, topsep=2pt]
    \item Cenário 1: $A_1 = 300,0 \times 0,40 \times 1,00 = 120,0\ \text{t}\cdot\text{ha}^{-1}\cdot\text{ano}^{-1}$ ($10$ vezes acima de $T = 12,0$);
    \item Cenário 2: $A_2 = 300,0 \times 0,05 \times 0,50 = 7,5\ \text{t}\cdot\text{ha}^{-1}\cdot\text{ano}^{-1}$ (abaixo de $T = 12,0$);
    \item Redução percentual: $\frac{120,0 - 7,5}{120,0} \times 100 = \frac{112,5}{120,0}\times 100 = 93,75\%$.
  \end{itemize}
  Além disso, a USLE foi calibrada para erosão laminar e em sulcos rasos, não simulando erosão subsuperficial (\textit{piping}) ou solapamento de taludes em voçorocas. \textit{Distratores:} (B) e (C) omitem $C_1$ ou $P_2$; (D) erra $A_1$ e atribui à USLE capacidade de modelar \textit{piping}; (E) inverte o efeito de $C$ e $P$.

  \item \textbf{Gabarito: (E).} Aplicando as fórmulas geométricas de dimensionamento de paliçadas em série:
  \begin{itemize}[leftmargin=1.4em, itemsep=1pt, topsep=2pt]
    \item Espaçamento horizontal: $E = \frac{h}{S} = \frac{0,60\ \text{m}}{0,10} = 6,00\ \text{m}$;
    \item Número de trechos/estruturas ao longo de $L = 60,0\ \text{m}$: $n = \frac{L}{E} = \frac{60,0}{6,00} = 10\ \text{paliçadas}$;
    \item Volume de sedimento retido por cunha: $V_{\text{cunha}} = \tfrac{1}{2}\,b\,h\,E = 0,5 \times 2,0 \times 0,60 \times 6,00 = 3,60\ \text{m}^3$;
    \item Capacidade total de retenção: $V_{\text{total}} = 10 \times 3,60 = 36,0\ \text{m}^3$.
  \end{itemize}
  \textit{Distratores:} (A) usa $S = 10$ em vez de $0,10$; (B) divide $E$ por $2$; (C) esquece o fator $\tfrac{1}{2}$ do prisma triangular ($b \cdot h \cdot E = 7,2\ \text{m}^3$ por estrutura); (D) dobra erroneamente o espaçamento para $12,0\ \text{m}$.

  \item \textbf{Gabarito: (C).} Pela Lei nº 12.651/2012:
  \begin{itemize}[leftmargin=1.4em, itemsep=1pt, topsep=2pt]
    \item \textbf{Art. 4º, I, \textit{b}:} para cursos d'água com largura de $10\ \text{m}$ a $50\ \text{m}$ (caso dos $35\ \text{m}$), a largura mínima da APP marginal é de \textbf{$50\ \text{m}$ em cada margem};
    \item \textbf{Área de APP marginal} (considerando as duas margens ao longo de $1.200\ \text{m}$):
    \[
      A_{\text{APP}} = 2 \times 50\ \text{m} \times 1.200\ \text{m} = 120.000\ \text{m}^2 = \frac{120.000}{10.000} = 12,0\ \text{ha};
    \]
    \item \textbf{Art. 12, II (Reserva Legal):} imóveis situados nas demais regiões do país fora da Amazônia Legal (incluindo o bioma Caatinga) devem manter \textbf{$20\%$} de Reserva Legal:
    \[
      A_{\text{RL}} = 0,20 \times 150,0\ \text{ha} = 30,0\ \text{ha}.
    \]
  \end{itemize}
  \textit{Distratores:} (A) aplica a faixa de rios $<10\ \text{m}$ ($30\ \text{m}$); (B) contabiliza apenas uma margem ($6,0\ \text{ha}$) e aplica $35\%$ (Cerrado na Amazônia Legal); (D) aplica a faixa de rios de $50$ a $200\ \text{m}$ ($100\ \text{m}$); (E) aplica $80\%$ (floresta na Amazônia Legal).

  \item \textbf{Gabarito: (D).} Um Diagnóstico Ambiental Aplicado e um PRAD tecnicamente defensáveis integram a escala da microbacia, removem a causa hidráulica da erosão antes da revegetação (drenagem e controle do escoamento a montante), empregam bioengenharia compatível com a intensidade do processo e definem indicadores verificáveis para monitoramento. \textit{Distratores:} (A) planta dentro de voçoroca ativa sem controle de fluxo (alta mortalidade das mudas); (B) carece de rigor quantitativo e cronograma; (C) impermeabiliza a bacia, agravando o escoamento e violando a função ecológica da APP; (E) substitui a recuperação obrigatória do passivo local por compensação externa inadequada.
\end{enumerate}
\end{boxGabarito}

% ============================================================================
\section*{Espelho de Correção das Questões Dissertativas (D1 a D5)}
% ============================================================================

\begin{boxConceito}[Critérios e Tópicos-Chave Esperados — Padrão Prova Presencial UEFS (0,8 pt cada)]
\small
\begin{enumerate}[leftmargin=1.6em, itemsep=5pt]
  \item \textbf{Dissertativa 1 (Módulo 1 — 0,8 pt):}
  \begin{itemize}[leftmargin=1.4em, itemsep=1pt, topsep=2pt]
    \item \textbf{Crítica epistemológica (0,3 pt):} explicitar que natureza e sociedade são indissociáveis; a crise ambiental decorre da racionalidade instrumental/econômica e da apropriação desigual dos recursos (Leff, 2001; Porto-Gonçalves, 2006), exigindo abordagem socioambiental (Mendonça, 2001).
    \item \textbf{Diferenciação conceitual (0,35 pt):} \textit{Perigo} = ameaça física potencial; \textit{Vulnerabilidade} = grau de exposição e fragilidade socioeconômica/estrutural de pessoas e bens; \textit{Risco} = combinação entre perigo e vulnerabilidade ($R = P \times V$); \textit{Resiliência} = capacidade do sistema de absorver perturbações e reorganizar suas funções.
    \item \textbf{Exemplificação contextualizada (0,15 pt):} caso concreto coerente na Bahia/Semiárido (ex.: secas prolongadas vs. vulnerabilidade hídrica rural; chuvas extremas em encostas urbanas ocupadas).
  \end{itemize}

  \item \textbf{Dissertativa 2 (Módulo 2 — 0,8 pt):}
  \begin{itemize}[leftmargin=1.4em, itemsep=1pt, topsep=2pt]
    \item \textbf{Etapas da AIA (0,25 pt):} triagem e escopo (Termo de Referência), diagnóstico, identificação/previsão/valoração de impactos, prognóstico, mitigação/compensação (PBA), decisão e monitoramento (Sánchez, 2013).
    \item \textbf{Diagnóstico vs. Prognóstico (0,3 pt):} diagnóstico retrata a linha de base atual dos meios físico, biótico e socioeconômico; prognóstico projeta e compara a evolução futura nos cenários \textit{com projeto} e \textit{sem projeto}.
    \item \textbf{Escalas e ordenamento (0,25 pt):} mostrar que fenômenos locais rápidos (erosão em parcela) acoplam-se a respostas regionais lentas (assoreamento de reservatórios), e que recortes espaciais/temporais subdimensionados mascaram impactos cumulativos e sinérgicos.
  \end{itemize}

  \item \textbf{Dissertativa 3 (Módulo 3 — 0,8 pt):}
  \begin{itemize}[leftmargin=1.4em, itemsep=1pt, topsep=2pt]
    \item \textbf{Bacia como unidade funcional (0,2 pt):} limites topográficos claros (divisores), convergência de água, sedimentos e solutos para um exutório comum e integração direta com o uso da terra (Botelho; Silva, 2004).
    \item \textbf{Morfopedologia de vertente (0,2 pt):} relação topo (infiltração/Latossolos) $\to$ meia encosta (escoamento/Neossolos e Cambissolos) $\to$ sopé (colúvios/Argissolos e Planossolos) $\to$ fundo de vale (aluviões/Gleissolos).
    \item \textbf{Tricart e Ross (0,4 pt):} articular meios estáveis (pedogênese > morfogênese), intergrades (equilíbrio dinâmico) e fortemente instáveis (morfogênese > pedogênese) à transição da Fragilidade Potencial (fatores físicos intrínsecos) para a Fragilidade Emergente (inserção do grau de proteção do uso e cobertura da terra).
  \end{itemize}

  \item \textbf{Dissertativa 4 (Módulo 4 — 0,8 pt):}
  \begin{itemize}[leftmargin=1.4em, itemsep=1pt, topsep=2pt]
    \item \textbf{Integração de indicadores do solo (0,3 pt):} justificar por que atributos físicos ($D_s$, $P_t$, infiltração, agregação), químicos ($CO/MO$, pH, CTC) e biológicos (atividade microbiana, fauna edáfica) avaliam de forma complementar as funções ecológicas do solo (Doran; Parkin, 1994; Vezzani; Mielniczuk, 2009).
    \item \textbf{Cadeia erosiva (0,25 pt):} perda de $MO$ e compactação reduzem estabilidade de agregados e infiltração $\to$ elevação do escoamento superficial $\to$ salpicamento (\textit{splash}) $\to$ erosão laminar $\to$ sulcos $\to$ ravinas $\to$ voçorocas (com eventual \textit{piping} ao atingir o lençol freático).
    \item \textbf{Aplicação do PEIR na Caatinga (0,25 pt):} propor ao menos um indicador mensurável coerente para Pressão, Estado, Impacto e Resposta.
  \end{itemize}

  \item \textbf{Dissertativa 5 (Módulo 5 — 0,8 pt):}
  \begin{itemize}[leftmargin=1.4em, itemsep=1pt, topsep=2pt]
    \item \textbf{Efeitos hidrológicos e mecânicos (0,35 pt):} hidrológicos (interceptação pela copa, serapilheira/rugosidade, infiltração e evapotranspiração que reduz poropressão $u$); mecânicos (reforço por tração de raízes finas elevando a coesão efetiva para $c' + c_r$ na equação $\tau_f = (c' + c_r) + (\sigma - u)\tan\phi'$, ancoragem profunda e contenção).
    \item \textbf{Bioengenharia, Código Florestal e Licenciamento (0,45 pt):} relacionar paliçadas, biomantas e estacas vivas à recomposição de APPs degradadas (Lei nº 12.651/2012) via PRAD e explicar sua vinculação às fases de LP (viabilidade e EIA/RIMA), LI (autorização de obras e PBA/medidas preventivas) e LO (operação condicionada ao monitoramento e cumprimento das condicionantes).
  \end{itemize}
\end{enumerate}
\end{boxConceito}

% ============================================================================
\section*{Painel de Diagnóstico Final por Módulo e Capítulo}
% ============================================================================

\noindent Utilize o painel abaixo para mapear seu desempenho por módulo e planejar a revisão direcionada para a Avaliação Teórica presencial e para a defesa do Trabalho Prático Final:

\begin{center}
\renewcommand{\arraystretch}{1.28}
\small
\begin{tabularx}{\textwidth}{@{}p{2.9cm} c c >{\raggedright\arraybackslash}X@{}}
  \toprule
  \rowcolor{uefsNavy}
  \textcolor{white}{\bfseries Módulo} & \textcolor{white}{\bfseries Questões} & \textcolor{white}{\bfseries Seu Acerto} & \textcolor{white}{\bfseries Rota de Revisão Prioritária (se $< 2/3$ obj. ou $< 0,5$ diss.)} \\
  \midrule
  \textbf{Mód. 1} \newline Fundamentos & Q01, Q02, Q03 \newline + D1 & \underline{\hspace{0.5cm}}/3 obj. \newline \underline{\hspace{0.5cm}}/0,8 & Revisar Caps. 1--4: epistemologia (Leff; Porto-Gonçalves), GTP (Bertrand), cadeia aspecto--impacto e matriz $R = P \times V$. \\
  \rowcolor{bgGray}
  \textbf{Mód. 2} \newline AIA e Integrada & Q04, Q05, Q06 \newline + D2 & \underline{\hspace{0.5cm}}/3 obj. \newline \underline{\hspace{0.5cm}}/0,8 & Revisar Caps. 5--6: etapas da AIA, hierarquia de mitigação, matriz ponderada $I_S$, retroalimentação e escalas. \\
  \textbf{Mód. 3} \newline Bacias e Ross & Q07, Q08, Q09 \newline + D3 & \underline{\hspace{0.5cm}}/3 obj. \newline \underline{\hspace{0.5cm}}/0,8 & Revisar Caps. 7--9: morfometria ($D_d, K_c, I_c, R_r$), declividade ($S\%, \theta^\circ$), Tricart e matriz de Ross. \\
  \rowcolor{bgGray}
  \textbf{Mód. 4} \newline Solo e Erosão & Q10, Q11, Q12 \newline + D4 & \underline{\hspace{0.5cm}}/3 obj. \newline \underline{\hspace{0.5cm}}/0,8 & Revisar Caps. 10--12: ciclo PEIR, cálculos de $D_s, P_t, MO$, normalização do IQS e sensibilidade da USLE. \\
  \textbf{Mód. 5} \newline Bioeng. e Leis & Q13, Q14, Q15 \newline + D5 & \underline{\hspace{0.5cm}}/3 obj. \newline \underline{\hspace{0.5cm}}/0,8 & Revisar Caps. 13--15: Mohr--Coulomb com $c_r$, paliçadas ($E = h/S$), APPs/RL (Lei 12.651/12), LP/LI/LO e PRAD. \\
  \midrule
  \multicolumn{4}{@{}l}{\textbf{Interpretação global:} \textbf{$\ge 8,0$ pts:} Domínio consolidado $\bullet$ \textbf{$6,0$ a $7,9$ pts:} Bom preparo (reforçar módulos abaixo de $70\%$) $\bullet$ \textbf{$< 6,0$ pts:} Retomar exemplos resolvidos (Etapa 4) e Formulário Técnico.} \\
  \bottomrule
\end{tabularx}
\end{center}
